%%%%%%%%%%  scan idx 504  |  folio 491  %%%%%%%%%%
\SGAPrintPage{491}%
Now by
definition we have:
\[
 T_i=-\cl(\underline{J}/\underline{J}^2)=-N;\qquad
 T_j=-\cl(\underline{J}'/\underline{J}'^2)=-L.
\]
Furthermore, the exact sequence:
\[
0\longrightarrow \underline{F}\longrightarrow g^*(\underline{N})\longrightarrow \underline{O}_{Y'}(1)\longrightarrow0
\]
gives the relation:
\[
 g^*(T_i)=-(F+L),
\]
from which we deduce:
\[
 \Todd(g^*(\check T_i))=\Todd(\check T_j)\cdot(\Todd(\check F))^{-1}.
\]

Taking 4.4 into account, the relation to be proved is therefore written:
\[
 \ch\bigl(j_*(g^*(x)\cdot\lambda_{-1}(F))\bigr)
 =
 j_{\mathrm{gr}}\bigl(\ch(g^*(x)\cdot\lambda_{-1}(F))\cdot\Todd(\check T_j)\bigr).
\]

Putting \(z=g^*(x)\cdot\lambda_{-1}(F)\), this relation is precisely the Riemann--Roch theorem applied to the regular immersion \(j\) of codimension \(1\), and to the element \(z\) of \(K^{\bullet}(Y')\).

Moreover, since we supposed that \(i\) satisfies the hypothesis of (VII 4.1), \(\lambda_{-1}(F)\) is divisible by \((1-L)\). One can therefore find \(z'\) such that \(z=z'(1-L)\), or equivalently \(z=j^*(j_*(z'))\) (exp. VII 2.7). We are thus reduced to the case where \(Y\) is a divisor, and to an element \(x\) such that there exists \(y\), with \(x=i^*(y)\).

\sgaStableTarget{sga6.e08.para.4-6}\textbf{4.6.} We must therefore show the relation:
\[
 \ch\bigl(i_*(i^*(y))\bigr)
 =
 i_{\mathrm{gr}}\bigl(\ch(i^*(y))\cdot\Todd(\check T_i)\bigr).
\]
Now the projection formula gives:
\[
 i_*(i^*(y))=y\cdot i_*(1),
\]
and from the exact sequence: \(0\longrightarrow \underline{J}\longrightarrow\underline{O}_X\longrightarrow\underline{O}_Y\longrightarrow0\), where \(\underline{J}\) is invertible here, we obtain:
\[
 i_*(1)=1-J,
\]
hence

%%%%%%%%%%  scan idx 505  |  folio 492  %%%%%%%%%%
\SGAPrintPage{492}%
\[
 \ch(i_*(i^*(y)))=\ch(y)\cdot\ch(1-J).
\]
Since \(\underline{J}\) is invertible, we have \(\widetilde c(J)=(1,1+c^1(J))\), hence \(\ch(J)=\exp(c^1(J))\).
Consequently, we have:
\[
 \ch(i_*(i^*(y)))=\ch(y)\cdot(1-\exp(c^1(J))).
\]

Furthermore, the functoriality of \(\ch\), and the projection formula 3.7 imply:
\[
 i_{\mathrm{gr}}\bigl(\ch(i^*(y))\cdot\Todd(\check T_i)\bigr)
 =
 \ch(y)\cdot i_{\mathrm{gr}}(\Todd(\check T_i)).
\]
Now \(T_i=-\cl(\underline{J}/\underline{J}^2)=-i^*(J)\). Thus:
\[
 i_{\mathrm{gr}}(\Todd(\check T_i))
 =
 i_{\mathrm{gr}}\bigl(i^{\mathrm{gr}}(\Todd(-\check J))\bigr)
 =
 i_{\mathrm{gr}}(1)\cdot\Todd(-\check J).
\]

Since \(i_*(1)=1-J\), \(i_{\mathrm{gr}}(1)\) is the image in \(\Gr^1K^{\bullet}(X)\) of \(1-J\), and is therefore \(-c^1(J)\). Furthermore, we have:
\[
 \Todd(-\check J)=(\Todd(\check J))^{-1}
 =
 \frac{1-\exp(c^1(J))}{-c^1(J)}.
\]

Thus finally we obtain:
\[
 i_{\mathrm{gr}}\bigl(\ch(i^*(y))\cdot\Todd(\check T_i)\bigr)
 =
 \ch(y)\cdot(1-\exp(c^1(J))),
\]
which proves the required relation.

\sgaStableTarget{sga6.e08.sec.5}\section*{5. Riemann--Roch theorem: case of the structural morphism of a projective bundle}

\sgaStableTarget{sga6.e08.para.5-1}\textbf{5.1.} By 3.9, to prove the Riemann--Roch theorem, it remains to prove it for the structural morphism \(f:\mathbf{P}^r_Y\longrightarrow Y\), where \(Y\) is a quasi-compact scheme having an ample sheaf. Put \(P=\mathbf{P}^r_Y\), \(\xi=\cl(\underline{O}_P(1))\), and \(\underline{E}=(\underline{O}_Y)^{r+1}\).

%%%%%%%%%%  scan idx 506  |  folio 493  %%%%%%%%%%
\SGAPrintPage{493}%
\sgaStableTarget{sga6.e08.prop.5-2}\textbf{Proposition 5.2.} \uline{We have on} \(P\) \uline{the exact sequence}:
\[
0\longrightarrow\underline{\Omega}^1_{P/Y}\longrightarrow f^*(\underline{E})(-1)\longrightarrow\underline{O}_P\longrightarrow0,
\]
\uline{i}.\uline{e}. \uline{we have an isomorphism}
\[
\underline{\Omega}^1_{P/Y}\simeq \underline{F}(-1),\quad\text{where}\quad \underline{F}=\Ker\bigl(f^*(\underline{E})\longrightarrow\underline{O}_P(1)\bigr).
\]

Consider the vector bundle over \(Y\) defined by \(\underline{E}\), and denote by \(Z\) the complement of the zero section of this bundle. We have the commutative diagram:
\[
\begin{tikzcd}[column sep=large,row sep=large]
P \arrow[dr,"f"'] && Z \arrow[ll,"h"'] \arrow[dl,"g"]\\
& Y. &
\end{tikzcd}
\]
The morphism \(h\) identifies \(Z\) with the affine \(P\)-scheme defined by the sheaf of algebras \(\underline{B}=\bigoplus_{n\in\mathbb Z}\underline{O}_P(n)\); \(\underline{B}\) is a graded \(\underline{O}_P\)-Algebra.

Since \(f,g,h\) are smooth morphisms, one may write the exact sequence of sheaves of relative differentials; since \(Z\) is affine over \(P\), one may write this sequence in terms of \(\underline{B}\)-Modules; one thus obtains:
\sgaStableTarget{sga6.e08.eq.5-2-1}\begin{equation}\tag{5.2.1}
0\longrightarrow
 \underline{\Omega}^1_{P/Y}\otimes_{\underline{O}_P}\underline{B}
\longrightarrow
 \underline{\Omega}^1_{\underline{B}/Y}
\longrightarrow
 \underline{\Omega}^1_{\underline{B}/\underline{O}_P}
\longrightarrow0 .
\end{equation}

Now this exact sequence is a graded exact sequence of graded \(\underline{B}\)-Modules. Indeed, it follows from the commutative diagram:
\[
\begin{tikzcd}[column sep=large,row sep=large]
\mathbf{S}(\underline{E}) \arrow[rr] && \displaystyle\bigoplus_{n\in\mathbb Z}\mathbf{S}(\underline{E})(n)\simeq \mathbf{S}(\underline{E})[T,T^{-1}]\\
& \underline{O}_Y \arrow[ul] \arrow[ur] &
\end{tikzcd}
\]
where the arrows are algebra homomorphisms compatible with the grading of \(\mathbf{S}(\underline{E})[T,T^{-1}]\) (grading defined by degree in the indeterminate \(T\)), the two other algebras being supposed concentrated

%%%%%%%%%%  scan idx 507  |  folio 494  %%%%%%%%%%
\SGAPrintPage{494}%
in degree \(0\).

The module of differentials \(\underline{\Omega}^1_{\underline{B}/\underline{O}_P}\) is isomorphic to \(\underline{B}(-1)\otimes_{\underline{O}_P}\underline{O}_P(1)\), since one can define an \(\mathbf{S}(\underline{E})\)-derivation \(d\):
\[
\mathbf{S}(\underline{E})[T,T^{-1}]
\longrightarrow
\mathbf{S}(\underline{E})[T,T^{-1}]\otimes_{\mathbf{S}(\underline{E})}\mathbf{S}(\underline{E})\cdot dT
\]
so as to satisfy the universal property defining \(\underline{\Omega}^1_{\underline{B}/\underline{O}_P}\), by setting \(d(aT^n)=naT^{n-1}\otimes dT\); moreover, \(d\) is compatible with the gradings if the factor \(\mathbf{S}(\underline{E})\cdot dT\) is assigned degree \(1\).

Let \(V\) be the vector bundle over \(Y\) defined by \(\underline{E}\). Then \(\underline{\Omega}^1_{Z/Y}\) is the restriction to \(Z\) of \(\underline{\Omega}^1_{V/Y}\). Now \(V\) is the affine \(Y\)-scheme defined by \(\mathbf{S}(\underline{E})\), so \(\underline{\Omega}^1_{V/Y}\) is defined by \(\mathbf{S}(\underline{E})\otimes_{\underline{O}_Y}\underline{E}\) (since \(\underline{E}\) is locally free of finite type), hence \(\underline{\Omega}^1_{V/Y}\) is equal to \(g^*(\underline{E})\), and consequently \(\underline{\Omega}^1_{\underline{B}/Y}=f^*(\underline{E})\otimes_{\underline{O}_P}\underline{B}\). To see its grading, one looks locally, and sees at once that \(f^*(\underline{E})\) must be regarded as being of degree \(1\).

Taking the degree \(0\) terms in the exact sequence (5.2.1) gives the announced exact sequence.

\sgaStableTarget{sga6.e08.rem.5-2-1}\textbf{Remark 5.2.1.} Using the general results on the differential study of Hilbert functors ([2], 5), one directly obtains an isomorphism \(\underline{\Omega}^1_{P/Y}\xrightarrow{\sim}\underline{F}(-1)\) from the given functorial characterization of \(g^*(\underline{\Omega}^1_{P/Y})\) for every \(Y\)-morphism \(g:Y'\to P\).

\sgaStableTarget{sga6.e08.para.5-3}\textbf{5.3.} By 5.2, we have:
\[
 T_f=\cl(\underline{\Omega}^1_{P/Y})=f^*(E)\cdot\xi^{-1}-1.
\]
Consequently:
\[
 \widetilde c(T_f)=\widetilde c(E)\cdot\widetilde c(\xi^{-1})-\widetilde c(1).
\]

%%%%%%%%%%  scan idx 508  |  folio 495  %%%%%%%%%%
\SGAPrintPage{495}%
Since \(\underline{E}=(\underline{O}_Y)^{r+1}\), we have \(E=r+1\), and consequently:
\[
 \widetilde c(E)=(r+1,1).
\]
Put:
\[
 x=\cl_{\Gr^1K^{\bullet}(P)}(\xi-1)
 =
 \cl_{\Gr^1K^{\bullet}(P)}(1-\xi^{-1})
 =
 c^1(\xi).
\]
Since \(\xi^{-1}=\check \xi\), we have:
\[
 \widetilde c(\xi^{-1})=(1,1-x).
\]
By V (6.2.1), this gives:
\[
 \widetilde c(T_f)=(r,(1-x)^{r+1}).
\]
Taking 3.5 into account, we finally obtain:
\[
 \widetilde c(\check T_f)=(r,(1+x)^{r+1}).
\]

By the definition of \(\Todd\), we therefore have:
\sgaStableTarget{sga6.e08.eq.5-3-1}\begin{equation}\tag{5.3.1}
 \Todd(\check T_f)=
 \left(\frac{x}{1-\exp(-x)}\right)^{r+1}.
\end{equation}

\sgaStableTarget{sga6.e08.para.5-4}\textbf{5.4.} Using the basis of \(K^{\bullet}(P)\) over \(K^{\bullet}(Y)\) formed by the elements \(\xi^k\), with \(-r\leq k\leq0\), one sees by additivity that it suffices to prove the theorem for elements of the form \(f^*(a)\cdot\xi^k\). We then have:
\[
 \ch\bigl(f_*(f^*(a)\xi^k)\bigr)
 =
 \ch(a f_*(\xi^k))
 =
 \ch(a)\cdot\ch(f_*(\xi^k)).
\]
On the other hand:
\[
 f_{\mathrm{gr}}\bigl(\ch(f^*(a)\xi^k)\Todd(\check T_f)\bigr)
 =
 f_{\mathrm{gr}}\bigl(f^{\mathrm{gr}}(\ch(a))\ch(\xi^k)\Todd(\check T_f)\bigr)
 =
 \ch(a)\cdot f_{\mathrm{gr}}\bigl(\ch(\xi^k)\Todd(\check T_f)\bigr).
\]

One thus sees that it suffices to prove the theorem for the \(\xi^k\), with \(-r\leq k\leq0\). Now by (VI (5.2.3)), we have
\[
 f_*(\xi^k)=0\quad\text{for }-r\leq k<0~;\qquad f_*(1)=1.
\]
Consequently:
\sgaStableTarget{sga6.e08.eq.5-4-1}\begin{equation}\tag{5.4.1}
 \ch(f_*(\xi^k))=0\quad\text{if }-r\leq k<0~;\qquad
 \ch(f_*(1))=1.
\end{equation}
